\scanpage{58}
\margin{\unsure{!}}
Un foncteur $E\leadsto\varphi_{T,a}(E)$ des $\pi_1(S,a)$-ensembles
finis dans lui-même, foncteur qui est
exact à gauche. Je dis que un tel
foncteur est nécessairement de la
forme $\underline{\operatorname{Hom}}_{\pi_1}(A,E)$, où $A$ est
un ens, dépendant \deleted{\unsure{tel quel|bien déf.}} fonctoriellement du
foncteur $\varphi_{T,\unsure{a|c}}$ qui lui donne naissance
[produit transformé en sommes].
% slanted
\margin{à détailler, et essayer de coiffer par un th. plus général. Il y a un \unsure{pble} \unsure{ou une proposition|une variante} à établir, qui résultera de la formule analogue à $(*)$ ci-dessous, où entre une $\varinjlim$.}
On \unsure{pose=|posera|pose} $A=A(T)$, \deleted{\unsure{la|On}} et on dénote
par $f_{\unsure{\ast|!}}^{\deleted{\unsure{\ast}}\mathrm{ét}}(T)$ la variété \deleted{\unsure{séparable|équivalent}} \inserted{étale} de $S$
associée. On a donc un isom de bifoncteurs
\[
\underline{\operatorname{Hom}}_S(f_\ast^{\mathrm{ét}}(T),S')=\deleted{\unsure{H}}\underline{\operatorname{Hom}}_S(T,S')
\tag*{\makebox[0pt][r]{\makebox[\linewidth][l]{$(*)$}}}
\]
[sur $\operatorname{Rev}\,\text{ét}\,(T)\times\operatorname{Rev}\,\text{ét}\,(S)$]

\deleted{d'autres termes pour \unsure{tt|tout} $T$, le foncteur
en $S'$ du deuxième mbr est
représentable}. \unsure{Prenant|Prenons} \deleted{\unsure{les sections|donc}} \inserted{$S'=T$}, \deleted{et
\unsure{faisant passer|Prenons}} les sections des deux
mbrs, on trouve
\[
\operatorname{Hom}_S(f_\ast^{\mathrm{ét}}(T),S')\simeq\operatorname{Hom}_S(T,S')
\]
d'où, en faisant $S'=f_\ast^{\unsure{\mathrm{ét}}}(T)$, \unsure{on trouve|un hom}
% archive sheet number: 57
